% Bagian: first-order-logic
% Bab: syntax-and-semantics
% Subbagian: terms-formulas

\documentclass[../../../include/open-logic-section]{subfiles}

\begin{document}

\olfileid[id]{fol}{syn}{frm}

\olsection{Suku dan \printtoken{P}{formula}}

Setelah suatu bahasa orde pertama~$\Lang L$ diberikan, kita dapat
mendefinisikan ungkapan yang dibangun dari kosakata dasar~$\Lang L$.
Ungkapan ini khususnya mencakup \emph{suku} dan \emph{!!{formula}s}.

\begin{defn}[Suku]
\ollabel{defn:terms}
Himpunan \emph{suku}~$\Trm[L]$ dari~$\Lang L$ didefinisikan secara
induktif sebagai berikut:
\begin{enumerate}
\item Setiap !!{variable} merupakan suku.
\item Setiap !!{constant} dari~$\Lang L$ merupakan suku.
\item Jika $f$ merupakan !!{function} $n$-tempat dan $t_1$, \dots, $t_n$
  merupakan suku, maka $\Atom{f}{t_1, \ldots, t_n}$ merupakan suku.
\tagitem{limitClause}{
  Tidak ada ungkapan lain yang merupakan suku.}{}
\end{enumerate}
Suku yang tidak memuat !!{variable} disebut \emph{suku tertutup}.
\end{defn}

\begin{explain}
Dalam spesifikasi bahasa dan suku, !!{constant}s muncul sebagai kategori
simbol tersendiri, tetapi kita juga dapat memasukkannya sebagai !!{function}s
nol-tempat. Dengan demikian, kita dapat menghilangkan klausa kedua dalam
definisi suku. Kita hanya perlu memahami $\Atom{f}{t_1, \ldots, t_n}$
sebagai $f$ itu sendiri jika $n = 0$.
\end{explain}

\begin{defn}[Formula]
\ollabel{defn:formulas}
Himpunan \emph{!!{formula}s}~$\Frm[L]$ dari bahasa~$\Lang L$
didefinisikan secara induktif sebagai berikut:
\begin{enumerate}
\tagitem{prvFalse}{$\lfalse$ merupakan !!{formula} atomik.}{}

\tagitem{prvTrue}{$\ltrue$ merupakan !!{formula} atomik.}{}

\item Jika $R$ merupakan !!{predicate} $n$-tempat dari~$\Lang L$ dan
  $t_1$, \dots, $t_n$ merupakan suku dari~$\Lang L$, maka
  $\Atom{R}{t_1,\ldots, t_n}$ merupakan !!{formula} atomik.

\item Jika $t_1$ dan $t_2$ merupakan suku dari~$\Lang L$, maka
  $\Atom{\eq}{t_1, t_2}$ merupakan !!{formula} atomik.

\tagitem{prvNot}{Jika $!A$ merupakan !!a{formula}, maka $\lnot !A$
  merupakan !!a{formula}.}{}

\tagitem{prvAnd}{Jika $!A$ dan $!B$ merupakan !!{formula}s, maka
  $(!A \land !B)$ merupakan !!a{formula}.}{}

\tagitem{prvOr}{Jika $!A$ dan $!B$ merupakan !!{formula}s, maka
  $(!A \lor !B)$ merupakan !!a{formula}.}{}

\tagitem{prvIf}{Jika $!A$ dan $!B$ merupakan !!{formula}s, maka
  $(!A \lif !B)$ merupakan !!a{formula}.}{}

\tagitem{prvIff}{Jika $!A$ dan $!B$ merupakan !!{formula}s, maka
  $(!A \liff !B)$ merupakan !!a{formula}.}{}

\tagitem{prvAll}{Jika $!A$ merupakan !!a{formula} dan $x$ merupakan
  !!a{variable}, maka $\lforall[x][!A]$ merupakan !!a{formula}.}{}

\tagitem{prvEx}{Jika $!A$ merupakan !!a{formula} dan $x$ merupakan
  !!a{variable}, maka $\lexists[x][!A]$ merupakan !!a{formula}.}{}

\tagitem{limitClause}{Tidak ada ungkapan lain yang merupakan
  !!a{formula}.}{}
\end{enumerate}
\end{defn}

\begin{explain}
Definisi himpunan suku dan himpunan !!{formula}s merupakan
\emph{definisi induktif}. Pada dasarnya, kita membangun himpunan
!!{formula}s dalam tak hingga banyaknya tahap. Pada tahap awal, kita
menetapkan semua formula atomik sebagai formula; hal ini bersesuaian dengan
beberapa kasus pertama dalam definisi, yaitu kasus untuk
\iftag{prvTrue}{$\ltrue$, }{}%
\iftag{prvFalse}{$\lfalse$, }{}%
$\Atom{R}{t_1,\dots,t_n}$ dan $\Atom{\eq}{t_1,t_2}$. Dengan demikian,
``!!{formula} atomik'' berarti setiap !!{formula} yang berbentuk demikian.

Kasus-kasus lain dalam definisi memberikan aturan untuk membangun
!!{formula}s baru dari !!{formula}s yang sudah dibangun. Pada tahap kedua,
kita dapat menggunakan aturan tersebut untuk membangun !!{formula}s dari
!!{formula}s atomik. Pada tahap ketiga, kita membangun formula baru dari
formula atomik dan formula yang diperoleh pada tahap kedua, dan demikian
seterusnya. Suatu ungkapan merupakan suatu !!{formula} jika pada akhirnya
dibangun pada salah satu tahap tersebut, dan tidak ada ungkapan lain yang
merupakan formula.
\end{explain}

Menurut konvensi, kita menulis $\eq$ di antara argumennya dan menghilangkan
tanda kurung: $\eq[t_1][t_2]$ merupakan singkatan untuk
$\Atom{\eq}{t_1,t_2}$. Selain itu, $\lnot \Atom{\eq}{t_1,t_2}$ disingkat
sebagai $\eq/[t_1][t_2]$. Ketika menulis formula $(!B \ast !C)$ yang
dibangun dari $!B$, $!C$ dengan menggunakan penghubung dua-tempat~$\ast$,
kita akan sering menghilangkan pasangan tanda kurung terluar dan cukup
menuliskan~$!B \ast !C$.

\begin{intro}
Beberapa teks logika mensyaratkan bahwa !!{variable}~$x$ harus muncul
dalam~$!A$ agar
\iftag{prvEx}{$\lexists[x][!A]$ }{}%
\iftag{notprvEx,notprvAll}{}{dan }%
\iftag{prvAll}{$\lforall[x][!A]$ }{}%
tergolong
\iftag{notprvEx,notprvAll}{!!a{formula}}{!!{formula}s}.
Tidak ada akibat buruk jika syarat ini tidak diberlakukan, dan hal itu
mempermudah pembahasan.
\end{intro}

\begin{tagblock}{defNot,defOr,defAnd,defIf,defIff,defTrue,defFalse,defEx,defAll}
\begin{defn}
Formula yang dibangun dengan menggunakan operator terdefinisi harus dipahami
sebagai berikut:

\begin{tagenumerate}{defTrue,defFalse,defNot,defOr,defAnd,defIf,defIff,defEx,defAll}

\tagitem{defTrue}{$\ltrue$ menyingkat
  \iftag{prvFalse}{$\lnot\lfalse$}{$(!A \lor \lnot !A)$ untuk suatu
    !!{formula} atomik tetap~$!A$}.}{}

\tagitem{defFalse}{$\lfalse$ menyingkat
  \iftag{prvTrue}{$\lnot\ltrue$}{$(!A \land \lnot !A)$ untuk suatu
    !!{formula} atomik tetap~$!A$}.}{}

\tagitem{defNot}{$\lnot !A$ menyingkat $!A \lif \lfalse$.}{}

\tagitem{defOr}{$!A \lor !B$ menyingkat
  \iftag{prvAnd}{$\lnot(\lnot !A \land \lnot !B)$}{$\lnot !A \lif
    !B$}.}{}

\tagitem{defAnd}{$!A \land !B$ menyingkat
  \iftag{prvOr}{$\lnot(\lnot !A \lor \lnot !B)$}{$\lnot (!A \lif
    \lnot !B)$}.}{}

\tagitem{defIf}{$!A \lif !B$ menyingkat
  \iftag{prvOr}{$\lnot !A \lor !B$}{$\lnot (!A \land \lnot !B)$}.}{}

\tagitem{defIff}{$!A \liff !B$ menyingkat $(!A \lif !B) \land (!B
  \lif !A)$.}{}

\tagitem{defAll}{$\lforall[x][!A]$ menyingkat
  $\lnot\lexists[x][\lnot !A]$.}{}

\tagitem{defEx}{$\lexists[x][!A]$ menyingkat
  $\lnot\lforall[x][\lnot !A]$.}{}
\end{tagenumerate}
\end{defn}
\end{tagblock}

Jika kita bekerja dalam suatu bahasa untuk penerapan tertentu, kita akan
sering menulis !!{predicate}s dan !!{function}s dua-tempat di antara
suku-suku yang bersangkutan, misalnya $t_1 < t_2$ dan $(t_1 + t_2)$ dalam
bahasa aritmetika serta $t_1 \in t_2$ dalam bahasa teori himpunan. Fungsi
penerus dalam bahasa aritmetika bahkan secara konvensional ditulis
\emph{setelah} argumennya:~$t'$. Namun, secara resmi, semuanya hanyalah
singkatan konvensional masing-masing untuk $\Atom{\Obj A^2_0}{t_1, t_2}$,
$\Obj f^2_0(t_1, t_2)$, $\Atom{\Obj A^2_0}{t_1, t_2}$, dan $f^1_0(t)$.

\begin{defn}[Identitas sintaktis]
Simbol $\ident$ menyatakan identitas sintaktis antara untai simbol, yakni
$!A \ident !B$ jika dan hanya jika $!A$ dan $!B$ merupakan untai simbol
dengan panjang yang sama dan memuat simbol yang sama pada setiap posisi.
\end{defn}

Simbol $\ident$ dapat diapit oleh untai yang diperoleh melalui konkatenasi;
misalnya, $!A \ident (!B \lor !C)$ berarti: untai simbol~$!A$ sama dengan
untai yang diperoleh dengan menggabungkan sebuah tanda kurung buka, untai
$!B$, simbol $\lor$, untai~$!C$, dan sebuah tanda kurung tutup, dalam urutan
tersebut. Jika demikian, kita mengetahui bahwa simbol pertama dari~$!A$
adalah tanda kurung buka, $!A$ memuat $!B$ sebagai subuntai (dimulai pada
simbol kedua), subuntai itu diikuti oleh~$\lor$, dan seterusnya.

Karena suku dan !!{formula}s dibangun dari unsur-unsur dasar melalui
definisi induktif, kita dapat menggunakan prinsip-prinsip induksi berikut
untuk membuktikan sifat-sifatnya.

\begin{lem}[\emph{Prinsip induksi pada suku}]
    \ollabel{lem:trmind}
    Misalkan $\Lang L$ suatu bahasa orde pertama.
    Jika suatu sifat~$P$ sedemikian rupa sehingga
    %
    \begin{enumerate}
        \item sifat itu berlaku bagi setiap !!{variable}~$v$,
        %
        \item sifat itu berlaku bagi setiap !!{constant}~$a$ dari~$\Lang L$,
          dan
        %
        \item sifat itu berlaku bagi $f(t_1,\dotsc,t_n)$ setiap kali sifat
          tersebut berlaku bagi $t_1$,~\dots, $t_n$ dan $f$ merupakan
          !!{function} $n$-tempat dari~$\Lang L$
    \end{enumerate}
    (dengan mengandaikan bahwa $t_1$,~\dots, $t_n$ merupakan suku
    dari~$\Lang{L}$), maka $P$ berlaku bagi setiap suku dalam~$\Trm[L]$.
\end{lem}

\begin{prob}
    Buktikan \olref[fol][syn][frm]{lem:trmind}.
\end{prob}

\begin{lem}[\emph{Prinsip induksi pada !!{formula}s}]
    \ollabel{thm:frmind}
    Misalkan $\Lang L$ suatu bahasa orde pertama.
    Jika suatu sifat~$P$ berlaku bagi semua !!{formula}s atomik dan
    sedemikian rupa sehingga
    %
    \begin{enumerate}
        \tagitem{prvNot}{sifat itu berlaku bagi $\lnot !A$ setiap kali
            berlaku bagi~$!A$;}{}
        \tagitem{prvAnd}{sifat itu berlaku bagi $(!A \land !B)$ setiap kali
            berlaku bagi $!A$ dan~$!B$;}{}
        \tagitem{prvOr}{sifat itu berlaku bagi $(!A \lor !B)$ setiap kali
            berlaku bagi $!A$ dan~$!B$;}{}
        \tagitem{prvIf}{sifat itu berlaku bagi $(!A \lif !B)$ setiap kali
            berlaku bagi $!A$ dan~$!B$;}{}
        \tagitem{prvIff}{sifat itu berlaku bagi $(!A \liff !B)$ setiap kali
            berlaku bagi $!A$ dan~$!B$;}{}
        \tagitem{prvEx}{sifat itu berlaku bagi $\lexists[x][!A]$ setiap kali
            berlaku bagi~$!A$;}{}
        \tagitem{prvAll}{sifat itu berlaku bagi $\lforall[x][!A]$ setiap kali
            berlaku bagi~$!A$;}{}
  \end{enumerate}
  (dengan mengandaikan bahwa $!A$ dan $!B$ merupakan !!{formula}s
  dari~$\Lang{L}$), maka $P$ berlaku bagi semua formula dalam~$\Frm[L]$.
\end{lem}

\end{document}

